\documentclass[11pt,a4paper]{article}
% =============================================================================
% Preambulo UTEQ - Aplicaciones Distribuidas
% Robusto contra overfull-hbox, colisiones de color y warnings de hyperref.
% Compatible con pdflatex + bibtex (ieeetr).
% =============================================================================

\usepackage[T1]{fontenc}
\usepackage[utf8]{inputenc}
\usepackage[spanish,es-tabla,es-nodecimaldot]{babel}
\usepackage{lmodern}
\usepackage{textcomp}
\usepackage{microtype}

% ---- Geometria conservadora (evita marginacion excesiva) --------------------
\usepackage[a4paper,margin=2.4cm,headheight=15pt,footskip=1cm]{geometry}

% ---- Manejo suave de justificacion (evita overfull hbox) --------------------
\emergencystretch=3em
\setlength{\parskip}{0.35em}
\setlength{\parindent}{0pt}

% ---- Colores institucionales UTEQ (nombres unicos, no colisionan) ----------
\usepackage[table,dvipsnames]{xcolor}
\definecolor{uteqBlue}{HTML}{0E3F7E}
\definecolor{uteqMid}{HTML}{1E5AA8}
\definecolor{uteqTeal}{HTML}{2E86A1}
\definecolor{uteqGreen}{HTML}{4FAE60}
\definecolor{uteqAmber}{HTML}{F2A413}
\definecolor{uteqGray}{HTML}{5A6470}
\definecolor{uteqLight}{HTML}{EEF2F8}
\definecolor{uteqSand}{HTML}{FAF3E6}
\definecolor{uteqRose}{HTML}{FBE9E7}

% ---- Empaquetados usuales ---------------------------------------------------
\usepackage{amsmath,amssymb,amsthm}
\usepackage{mathtools}
\usepackage{graphicx}
\usepackage{booktabs}
\usepackage{array}
\usepackage{makecell}
\usepackage{longtable}
\usepackage{xltabular}
\usepackage{ragged2e}
\usepackage{enumitem}
\setlist{itemsep=2pt,topsep=3pt,parsep=0pt}
\usepackage{fancyhdr}
\usepackage{titlesec}
\usepackage{caption}
\usepackage[section]{placeins}

\captionsetup[table]{labelfont=bf,font=small,skip=4pt}
\captionsetup[figure]{labelfont=bf,font=small,skip=4pt}

% ---- Encabezado y pie -------------------------------------------------------
\pagestyle{fancy}
\fancyhf{}
\renewcommand{\headrulewidth}{0.4pt}
\renewcommand{\footrulewidth}{0pt}
\fancyhead[L]{\small\textcolor{uteqBlue}{UTEQ $\cdot$ Ingeniería de Software}}
\fancyhead[R]{\small\textcolor{uteqGray}{Aplicaciones Distribuidas $\cdot$ ISR-701}}
\fancyfoot[C]{\thepage}

% ---- Titulos con color institucional ---------------------------------------
\titleformat{\section}{\normalfont\Large\bfseries\color{uteqBlue}}{\thesection.}{0.6em}{}
\titleformat{\subsection}{\normalfont\large\bfseries\color{uteqMid}}{\thesubsection.}{0.5em}{}
\titleformat{\subsubsection}{\normalfont\normalsize\bfseries\color{uteqTeal}}{\thesubsubsection.}{0.5em}{}
\titlespacing*{\section}{0pt}{1.2em}{0.6em}
\titlespacing*{\subsection}{0pt}{0.8em}{0.4em}
\titlespacing*{\subsubsection}{0pt}{0.6em}{0.3em}

% ---- Cuadros de texto -------------------------------------------------------
\usepackage[most]{tcolorbox}
\newtcolorbox{utqbox}[1][]{
  enhanced,
  colback=uteqLight,
  colframe=uteqBlue,
  boxrule=0.7pt,
  arc=1.2mm,
  left=6pt,right=6pt,top=4pt,bottom=4pt,
  fonttitle=\bfseries\color{white},
  coltitle=white,
  colbacktitle=uteqBlue,
  attach boxed title to top left={xshift=6pt,yshift=-2mm},
  boxed title style={arc=1mm,colback=uteqBlue},
  #1
}
\newtcolorbox{noteBox}[1][]{
  enhanced,colback=uteqSand,colframe=uteqAmber,boxrule=0.6pt,arc=1mm,
  left=6pt,right=6pt,top=3pt,bottom=3pt,#1
}
\newtcolorbox{critBox}[1][]{
  enhanced,colback=uteqRose,colframe=uteqBlue,boxrule=0.6pt,arc=1mm,
  left=6pt,right=6pt,top=3pt,bottom=3pt,#1
}

% ---- Codigo fuente (lstlisting) --------------------------------------------
\usepackage{listings}

% Definicion manual del lenguaje YAML para listings
\lstdefinelanguage{yaml}{
  keywords={true,false,null,yes,no,on,off},
  keywordstyle=\color{uteqBlue}\bfseries,
  sensitive=false,
  comment=[l]{\#},
  morecomment=[s]{---}{---},
  commentstyle=\color{uteqGreen}\itshape,
  stringstyle=\color{uteqAmber},
  moredelim=[l][\color{uteqTeal}]{:},
  morestring=[b]',
  morestring=[b]"
}

\lstdefinestyle{utqcode}{
  basicstyle=\ttfamily\scriptsize,
  keywordstyle=\color{uteqBlue}\bfseries,
  commentstyle=\color{uteqGreen}\itshape,
  stringstyle=\color{uteqAmber},
  numbers=left,
  numberstyle=\tiny\color{uteqGray},
  numbersep=6pt,
  frame=single,
  rulecolor=\color{uteqGray},
  backgroundcolor=\color{uteqLight!45},
  breaklines=true,
  breakatwhitespace=true,
  breakindent=8pt,
  postbreak=\mbox{\textcolor{uteqAmber}{$\hookrightarrow$}\space},
  columns=fullflexible,
  keepspaces=true,
  showstringspaces=false,
  captionpos=b,
  aboveskip=8pt,belowskip=8pt,
  xleftmargin=10pt,xrightmargin=4pt,
  tabsize=2,
  extendedchars=true,
  inputencoding=utf8,
  literate=%
    {á}{{\'a}}1 {é}{{\'e}}1 {í}{{\'i}}1 {ó}{{\'o}}1 {ú}{{\'u}}1
    {Á}{{\'A}}1 {É}{{\'E}}1 {Í}{{\'I}}1 {Ó}{{\'O}}1 {Ú}{{\'U}}1
    {ñ}{{\~n}}1 {Ñ}{{\~N}}1 {ü}{{\"u}}1 {Ü}{{\"U}}1
    {¿}{{?`}}1 {¡}{{!`}}1
}
\lstset{style=utqcode}

% ---- Enlaces (hyperref sin warnings de Unicode) ----------------------------
\usepackage[hidelinks,pdfusetitle,unicode=true]{hyperref}
\usepackage{url}

% ---- Simbolos de check / cruz ----------------------------------------------
\usepackage{pifont}
\newcommand{\ok}{\textcolor{uteqGreen}{\ding{51}}}
\newcommand{\ko}{\textcolor{red!70!black}{\ding{55}}}

% ---- Cabecera institucional reutilizable -----------------------------------
\newcommand{\uteqheader}[3]{%
  \begin{utqbox}[title={Universidad Técnica Estatal de Quevedo $\cdot$ Facultad de Ciencias de la Computación}]
    \small
    \textbf{Carrera:} Ingeniería de Software \hfill \textbf{Asignatura:} Aplicaciones Distribuidas (ISR-701)\\
    \textbf{Documento:} #1 \hfill \textbf{Semana de entrega:} #2\\
    \textbf{Período:} Regular 2026-2027 \hfill \textbf{Modalidad:} Grupal (equipos PFC de 3-4 integrantes)\\
    \textbf{Docente responsable:} Ing.\ Gleiston C.\ Guerrero-Ulloa, Mgs.\ \hfill \textbf{Versión:} #3
  \end{utqbox}
}


\title{Guia de la Entrega 4 del Proyecto Fin de Curso}
\author{Aplicaciones Distribuidas $\cdot$ ISR-701}
\date{}

\begin{document}
\thispagestyle{fancy}

\begin{center}
    {\Large\bfseries\color{uteqBlue} Guia Integral de la Entrega 4 del Proyecto Fin de Curso}\\[2pt]
    {\normalsize\color{uteqGray} Aplicación completa (Web y Móvil), calidad y evaluación experimental}\\[2pt]
    {\small Aplicaciones Distribuidas $\cdot$ Septimo semestre $\cdot$ ISR-701}
\end{center}
\vspace{0.4em}

\uteqheader{Guia PFC Entrega 4 (E4) --- Entrega final}{Semana 17}{v1.0}

\section{Presentación y lógica acumulativa de la entrega}

    La Entrega 4 (E4) constituye el cierre integral del PFC. En ella el sistema distribuido, construido incrementalmente durante todo el semestre, se completa con sus dos interfaces obligatorias (aplicación web y aplicación móvil), su capa de calidad (pruebas automatizadas y \textit{pipeline} de integración y despliegue continuos), su infraestructura de observabilidad (métricas, registros y trazas distribuidas) y su evaluación experimental frente al modelo de calidad ISO/IEC 25010:2023~\cite{ISO25010_2023}.
    
    La E4 hereda directamente los artefactos producidos en las entregas anteriores: el análisis del dominio y el prototipo de comunicación de la E1, la arquitectura de microservicios REST con Docker Compose de la E2, y la capa de datos distribuida en CockroachDB con el pipeline paralelo PySpark de la E3. Ninguna de esas piezas se desecha. La E4 las envuelve, las refactoriza para elevar su cohesion y su calidad, y las expone a los usuarios finales a través de dos canales: un cliente web y un cliente móvil.
    
    \begin{critBox}
        \textbf{Requisito estricto de esta entrega.} Cada equipo debe presentar dos aplicaciones cliente que consumen la misma API distribuida: (i) una aplicación web y (ii) una aplicación móvil. Ambas son obligatorias y no sustituibles entre si. La ausencia total de una de ellas aplica un factor multiplicativo de $0{,}5$ sobre la nota final de la entrega.
    \end{critBox}
    
    El nivel de rigor exigido corresponde al reporte de un experimento controlado en ingeniería de software: propósito, hipótesis, variables, protocolo, replicabilidad y discusión de amenazas a la validez~\cite{Wohlin2012Experimentation,Runeson2009CaseStudy}. El manuscrito final consolidara todo el trabajo del semestre y será la evidencia única de evaluación del cierre.

\section{Objetivos}

\subsection{Objetivo general}

    Completar, evaluar y documentar el sistema distribuido del PFC con arquitectura en capas, aplicaciones web y móvil, pipeline CI/CD, observabilidad distribuida y evaluación experimental de calidad, cumpliendo el modelo ISO/IEC 25010:2023 y los principios de reproducibilidad propios de la investigación empírica.

\subsection{Objetivos específicos}

    \begin{enumerate}
        \item Refactorizar el backend de E2--E3 aplicando arquitectura en capas, principios SOLID~\cite{Martin2017CleanArch} y al menos cinco patrones de diseño GoF~\cite{Gamma1994GoF}.
        \item Desarrollar una aplicación web (SPA) que consume la API a través del API Gateway, con \textit{routing}, gestión de estado, internacionalizacion y accesibilidad básica.
        \item Desarrollar una aplicación móvil (Android nativa con Kotlin o multiplataforma con Flutter/React Native) que consume la misma API y hace uso de al menos dos capacidades del dispositivo (sensores, camara, notificaciones o geolocalizacion).
        \item Establecer un \textit{pipeline} completo de CI/CD con GitHub Actions que compila, prueba, empaqueta y pública artefactos de las tres capas (backend, web, móvil)~\cite{Kim2021DevOpsHandbook,Humble2010CD}.
        \item Instrumentar el sistema con las tres señales de observabilidad (métricas, registros estructurados y trazas distribuidas) según la especificacion OpenTelemetry~\cite{OpenTelemetry2024,Sigelman2010Dapper} y publicar un dashboard operativo.
        \item Evaluar el sistema contra cinco características del modelo ISO/IEC 25010:2023 con métricas objetivas medidas experimentalmente~\cite{ISO25010_2023}.
        \item Presentar el manuscrito final consolidando la totalidad del ciclo del PFC, con al menos 20 páginas de contenido, con estructura de reporte experimental y con análisis de amenazas a la validez.
    \end{enumerate}

\section{Aplicación por equipo PFC}

    La Tabla~\ref{tab:pfc-e4-app} concreta la aplicación por equipo. Se supone conocida la arquitectura de microservicios de E2 y la capa de datos distribuida de E3; la E4 se centra en refactorizar el backend en capas, materializar los dos clientes y evaluar la calidad.

    \renewcommand{\arraystretch}{1.15}
    \small
    \begin{xltabular}{\linewidth}{@{}p{1.9cm} X X X@{}}
        \caption{Aplicación concreta de E4 por equipo PFC.}
        \label{tab:pfc-e4-app}\\
        \toprule
        \textbf{Equipo} & \textbf{Aplicación Web (SPA)} & \textbf{Aplicación Móvil (obligatoria)} & \textbf{Patrones GoF aplicados}\\
        \midrule
        \endfirsthead
        \toprule
        \textbf{Equipo} & \textbf{Aplicación Web (SPA)} & \textbf{Aplicación Móvil (obligatoria)} & \textbf{Patrones GoF aplicados}\\
        \midrule
        \endhead
        \bottomrule
        \endfoot
        AGLS \newline TiendaTech &
        Catalogo, carrito, checkout, panel de vendedor, dashboard de métricas para el administrador. &
        Cliente móvil de compra: catalogo, carrito, notificaciones \textit{push} de estado del pedido, camara para reconocer códigos de barras. &
        Repository, Factory Method, Strategy (métodos de pago), Observer (eventos de pedido), Decorator (promociones).\\
        \addlinespace
        BCEL \newline SGA Escuela &
        Portal de docente, secretaría y directivo, con captura de calificaciones y reportes por período. &
        Cliente móvil para representantes: consulta de calificaciones, asistencia y comunicaciones, con notificaciones \textit{push}. &
        Repository, Template Method (reportes), Observer (notificaciones), Facade (fachada del portal), Strategy (algoritmos de promedio).\\
        \addlinespace
        FUVV \newline Laboratorios &
        Panel de reserva y monitoreo con visualizacion de series temporales y calendario semanal. &
        Cliente móvil para técnicos: reservas rapidas, notificacion de incidentes, escaneo de QR de sala. &
        Singleton (cliente Prometheus), Observer (alertas), State (estados de la reserva), Strategy (políticas de reserva), Facade.\\
        \addlinespace
        ACC \newline Soporte ISP &
        Consola de operadores y coordinadores, con vista de SLA por zona y por técnico. &
        Cliente móvil para técnicos de campo: tickets asignados, cierre en sitio, geolocalizacion y camara para evidencias. &
        Command (acciones sobre tickets), Chain of Responsibility (escalado), Observer (SLA), Repository, Strategy (rutas).\\
    \end{xltabular}
    \normalsize

\section{Fundamento teórico exigible}

\subsection{Arquitectura en capas y principios SOLID}

    La arquitectura en capas separa el sistema en (i) presentación, (ii) aplicación, (iii) dominio e (iv) infraestructura~\cite{Martin2017CleanArch,Evans2003DDD}. Los principios SOLID son las restricciones que impiden que las capas se acoplen incorrectamente y que las clases se conviertan en \emph{objetos-Dios}: responsabilidad única, abierto-cerrado, sustitución de Liskov, segregacion de interfaz e inversion de dependencias. El uso disciplinado de estos principios es prerrequisito para que los patrones GoF~\cite{Gamma1994GoF} aporten valor real.

\subsection{Ingeniería de aplicaciones móviles}

    Las aplicaciones móviles se rigen por retos particulares: heterogeneidad de dispositivos, restricciones de batería y red, ciclo de vida gestionado por el sistema operativo y ventanas de publicación controladas por tiendas~\cite{Wasserman2010MobileSE}. La decisión entre desarrollo nativo o multiplataforma no es una preferencia estetica: tiene consecuencias medibles en rendimiento, mantenimiento y \textit{time-to-market}~\cite{ElKassas2017CrossPlatform,BiornHansen2019CrossPlatform,Ahmad2018CrossPlatformSurvey}. El equipo debe justificar su elección con al menos dos criterios cuantitativos.

\subsection{Integración y entrega continuas (CI/CD)}

    La entrega continua se apoya en un \textit{pipeline} automatizado que ejecuta pruebas, empaqueta artefactos y despliega en ambientes segregados~\cite{Humble2010CD}. La cultura DevOps subyacente~\cite{Kim2021DevOpsHandbook} demanda que ningun cambio llegue a producción sin pasar por los mismos controles. En el PFC el \emph{pipeline} de GitHub Actions actua como \emph{quality gate}: si la integración falla, el despliegue no ocurre.

\subsection{Observabilidad distribuida}

    La observabilidad de un sistema distribuido se sustenta en tres señales: métricas (agregados numéricos con etiquetas), registros (\textit{logs} estructurados) y trazas (secuencias de \textit{spans} correlacionados por un identificador comun)~\cite{Sigelman2010Dapper,Beyer2016SRE}. El estándar OpenTelemetry~\cite{OpenTelemetry2024} unifica su modelo y su exportacion, permitiendo que diferentes proveedores consuman las mismas señales. En el PFC estas tres señales se integran en un dashboard operativo que se demostrara en vivo.

\subsection{Modelo de calidad ISO/IEC 25010:2023}

    El modelo ISO/IEC 25010:2023~\cite{ISO25010_2023} organiza la calidad de producto en ocho características: adecuación funcional, eficiencia de desempeño, compatibilidad, usabilidad, fiabilidad, seguridad, mantenibilidad y portabilidad. Cada característica se descompone en subcaracterísticas y admite medición mediante métricas objetivas. En el PFC se evaluarán cinco de ellas con al menos una métrica cuantitativa cada una.

\section{Manual de desarrollo (paso a paso)}

    Se asume que el equipo parte de la rama \texttt{main} con la E3 fusionada. Todo el trabajo de esta entrega se realiza en la rama \texttt{feature/entrega-4} y se fusiona a \texttt{main} por \textit{pull request} tras la revisión cruzada final.

\subsection{Stack tecnológico consolidado}

    \begin{utqbox}[title={Stack obligatorio para E4}]
        \small
        \begin{itemize}
            \item \textbf{Backend:} Java 21 + Spring Boot 3.x (o el equivalente adoptado en E2), refactorizado en capas.
            \item \textbf{Persistencia distribuida:} cluster CockroachDB de la E3, sin cambios.
            \item \textbf{Analitica paralela:} pipeline PySpark de la E3, empaquetado en el mismo \textit{compose}.
            \item \textbf{Web (SPA):} React 18 o Angular 17 (a elección del equipo, justificada); TypeScript; \textit{router} declarativo; \textit{state} con Redux Toolkit, Zustand, NgRx o Signals; \textit{i18n} con \texttt{react-i18next} o \texttt{ngx-translate}; empaquetado en Docker \textit{multi-stage} servido con Nginx.
            \item \textbf{Móvil:} Android nativa con Kotlin y Jetpack Compose, o multiplataforma con Flutter 3.x o React Native 0.74+.
            \item \textbf{CI/CD:} GitHub Actions con \textit{matrix build} para las tres capas.
            \item \textbf{Observabilidad:} Prometheus (heredado de E2), Grafana, OpenTelemetry Collector, exportadores JVM y de sistema.
            \item \textbf{Pruebas de carga:} Locust 2.x.
            \item \textbf{Documentación:} LaTeX (repositorio ya lo tiene desde E1).
        \end{itemize}
    \end{utqbox}

\subsection{Módulo A: refactorización del backend en capas}

    \textbf{Entrada:} microservicio principal de E3 conectado al cluster CockroachDB.\\
    \textbf{Procedimiento:}
    \begin{enumerate}
        \item Reorganizar el código en cuatro paquetes por microservicio: \path{presentation}, \path{application}, \path{domain} e \path{infrastructure}.
        \item Extraer entidades y objetos de valor al paquete \path{domain}; en Spring Boot, estas entidades \emph{no} deben tener anotaciones JPA (las anotaciones se colocan en clases de infraestructura).
        \item Introducir puertos (interfaces) en \path{domain} y adaptadores en \path{infrastructure}: repositorios JDBC, clientes HTTP y \textit{publishers} de mensajes.
        \item Aplicar cinco patrones GoF con nombre y propósito documentados en un ADR (\path{docs/adr/ADR-005-patrones-gof.md}). La lista puede variar por equipo pero debe incluir al menos Repository, Factory Method y Strategy.
        \item Actualizar las pruebas de E2 y E3 para que se apoyen en dobles de prueba de los puertos y no en la infraestructura real, exceptuando las pruebas de integración.
    \end{enumerate}
    \textbf{Salida esperada:} el sistema mantiene el mismo comportamiento externo (API sin cambios), pero su estructura interna esta desacoplada. La cobertura de pruebas unitarias sube a $\ge 70\%$.

\subsection{Módulo B: aplicación web (SPA)}

    \textbf{Procedimiento (independiente del framework elegido):}
    \begin{enumerate}
        \item Crear la carpeta \path{apps/web/} en el repositorio.
        \item Configurar el proyecto con TypeScript en modo estricto y \textit{linter} obligatorio (\texttt{eslint} + \texttt{prettier}).
        \item Implementar el patrón de arquitectura por características (\textit{feature-based}): cada módulo de negocio agrupa componentes, servicios de acceso a la API, hooks (o servicios) y pruebas.
        \item Cubrir al menos cinco rutas: \path{/}, \path{/login}, \path{/main} (el dominio del PFC), \path{/settings} y \path{/about}; incluir manejo de rutas protegidas por JWT.
        \item Integrar internacionalizacion (español e ingles como mínimo), tema claro/oscuro y estados de carga y de error explicitos.
        \item Consumir la API del backend a través del API Gateway; guardar el JWT en \texttt{httpOnly cookie} o en \texttt{sessionStorage} con expiracion, nunca en \texttt{localStorage} plano.
        \item Empaquetar en \textit{Dockerfile multi-stage} (etapa \texttt{builder} con Node LTS y etapa \texttt{runtime} con \texttt{nginx:alpine}); agregar el servicio al \texttt{docker-compose.yml}.
        \item Publicar como \emph{artefacto} en el pipeline CI un \texttt{.tar.gz} del \path{dist/} versionado con el hash de \textit{commit}.
    \end{enumerate}
    El Listado~\ref{lst:dockerfile-web} muestra el Dockerfile de referencia.
    \begin{lstlisting}[language=bash,caption={Dockerfile multi-stage para la SPA.},label={lst:dockerfile-web}]
# apps/web/Dockerfile
FROM node:20-alpine AS builder
WORKDIR /app
COPY package*.json ./
RUN npm ci
COPY . .
RUN npm run build

FROM nginx:1.27-alpine
COPY --from=builder /app/dist /usr/share/nginx/html
COPY nginx.conf /etc/nginx/conf.d/default.conf
EXPOSE 80
HEALTHCHECK --interval=30s --timeout=3s \
  CMD wget -q -O- http://localhost/ || exit 1
    \end{lstlisting}

\subsection{Módulo C: aplicación móvil}

    \textbf{Procedimiento (Android nativa como referencia; los equipos que elijan multiplataforma siguen los pasos analogos):}
    \begin{enumerate}
        \item Crear \path{apps/mobile/} con proyecto Gradle Android; SDK mínimo API 26 (Android 8), objetivo API 34.
        \item Definir la arquitectura MVVM con Jetpack Compose para la vista, ViewModel para el estado y un repositorio por característica.
        \item Implementar la autenticacion consumiendo el mismo endpoint \path{/auth/login} del backend; almacenar el JWT en \texttt{EncryptedSharedPreferences}.
        \item Implementar el listado y el detalle del recurso principal del PFC, con \textit{pull to refresh} y modo sin conexión (caché local con Room o SQLDelight).
        \item Integrar al menos dos capacidades del dispositivo: notificaciones \textit{push} (Firebase Cloud Messaging), camara (para OCR o códigos), geolocalizacion (fused location), o sensores de aceleración según el dominio del PFC.
        \item Cubrir con pruebas unitarias los ViewModels (JUnit 5 + coroutines-test) y con al menos una prueba instrumentada \textit{end-to-end} (Espresso o Compose Testing).
        \item Generar un APK firmado con clave de \emph{debug} y publicarlo como artefacto del pipeline CI en \path{release/apk/}.
    \end{enumerate}

\subsection{Módulo D: pirámide de pruebas}

    La estrategia de pruebas cubre las cuatro capas de la pirámide clasica~\cite{Ammann2016SWTesting,Winters2020SEatGoogle}:
    \begin{itemize}
        \item \textbf{Unitarias:} cobertura $\ge 70\%$ en backend, web y móvil. Herramientas: JUnit 5 / Vitest o Jest / JUnit Android.
        \item \textbf{Integración:} pruebas de repositorios con Testcontainers apuntando a un CockroachDB efímero.
        \item \textbf{Contrato:} \emph{consumer-driven contract testing} con Pact entre el frontend (web y móvil como consumidores) y el backend (proveedor).
        \item \textbf{Extremo a extremo (E2E):} Playwright para la web; Espresso o Detox para el móvil.
        \item \textbf{Carga:} Locust con al menos dos escenarios: (i) 50 usuarios durante 5 minutos, (ii) rampa progresiva de 0 a 200 usuarios en 10 minutos.
    \end{itemize}
    Cada tipo de prueba se ejecuta en un \textit{job} distinto del pipeline y sus resultados se almacenan como artefactos.

\subsection{Módulo E: pipeline CI/CD end-to-end}

    El archivo \path{.github/workflows/ci-cd.yml} define siete \textit{jobs} en un \emph{DAG}: \texttt{lint}, \texttt{test-backend}, \texttt{test-web}, \texttt{test-mobile}, \texttt{build-images}, \texttt{build-mobile-apk} e \texttt{integration}. Solo cuando todos los \textit{jobs} estan en verde, el pipeline pública las imagenes en GitHub Container Registry (GHCR) etiquetadas con el hash del commit. El Listado~\ref{lst:ci} muestra un esqueleto reducido del pipeline.
    \begin{lstlisting}[language=yaml,caption={Esqueleto del pipeline .github/workflows/ci-cd.yml.},label={lst:ci}]
name: pfc-cicd
on:
  push:    { branches: [main, "feature/*"] }
  pull_request:

jobs:
  lint:
    runs-on: ubuntu-latest
    steps:
      - uses: actions/checkout@v4
      - name: Backend lint
        run: ./mvnw -pl services/resource-service checkstyle:check
      - name: Web lint
        working-directory: apps/web
        run: npm ci && npm run lint

  test-backend:
    needs: lint
    runs-on: ubuntu-latest
    steps:
      - uses: actions/checkout@v4
      - uses: actions/setup-java@v4
        with: { distribution: temurin, java-version: "21" }
      - run: ./mvnw -pl services/resource-service verify

  test-web:
    needs: lint
    runs-on: ubuntu-latest
    steps:
      - uses: actions/checkout@v4
      - working-directory: apps/web
        run: |
          npm ci
          npm run test -- --coverage

  test-mobile:
    needs: lint
    runs-on: ubuntu-latest
    steps:
      - uses: actions/checkout@v4
      - uses: actions/setup-java@v4
        with: { distribution: temurin, java-version: "17" }
      - working-directory: apps/mobile
        run: ./gradlew test

  build-images:
    needs: [test-backend, test-web]
    runs-on: ubuntu-latest
    permissions: { contents: read, packages: write }
    steps:
      - uses: actions/checkout@v4
      - name: Login GHCR
        uses: docker/login-action@v3
        with:
          registry: ghcr.io
          username: ${{ github.actor }}
          password: ${{ secrets.GITHUB_TOKEN }}
      - name: Build backend
        run: docker build -t ghcr.io/${{ github.repository }}/backend:${{ github.sha }} services/resource-service
      - name: Build web
        run: docker build -t ghcr.io/${{ github.repository }}/web:${{ github.sha }} apps/web
      - run: docker push --all-tags ghcr.io/${{ github.repository }}/backend
      - run: docker push --all-tags ghcr.io/${{ github.repository }}/web

  build-mobile-apk:
    needs: test-mobile
    runs-on: ubuntu-latest
    steps:
      - uses: actions/checkout@v4
      - working-directory: apps/mobile
        run: ./gradlew assembleDebug
      - uses: actions/upload-artifact@v4
        with:
          name: pfc-debug-apk
          path: apps/mobile/app/build/outputs/apk/debug/*.apk

  integration:
    needs: build-images
    runs-on: ubuntu-latest
    steps:
      - uses: actions/checkout@v4
      - run: docker compose -f docker-compose.yml up -d
      - run: |
          sleep 30
          curl -f http://localhost:8080/actuator/health
      - run: docker compose down
    \end{lstlisting}

\subsection{Módulo F: observabilidad distribuida}

    \textbf{Procedimiento:}
    \begin{enumerate}
        \item Añadir al microservicio principal el \emph{starter} \texttt{opentelemetry-spring-boot-starter} (o equivalente).
        \item Configurar el exportador OTLP hacia un \textit{OpenTelemetry Collector} añadido al \texttt{docker-compose.yml}.
        \item Reutilizar las métricas Prometheus de la E3 (\path{crdb_*}) y añadir cuatro nuevas: \path{http_requests_total} (etiquetado por ruta, método y código), \path{http_request_duration_seconds}, \path{app_business_events_total} y \path{app_active_sessions}.
        \item Emitir registros en formato JSON con Logback (backend) y con \texttt{Timber} + \texttt{tree} JSON (móvil), incluyendo el campo \texttt{trace\_id} para correlacion.
        \item Construir un dashboard Grafana con seis paneles obligatorios: tasa de peticiones, latencia p50/p95/p99, tasa de errores 5xx, disponibilidad Raft del cluster CRDB, uso de recursos por contenedor y latencia \textit{end-to-end} desde el cliente móvil.
        \item Exportar el dashboard como JSON versionado en \path{ops/grafana/pfc-dashboard.json} para que sea reproducible.
    \end{enumerate}

\subsection{Módulo G: evaluación experimental ISO 25010}

    \textbf{Protocolo experimental (rigor JCR):}
    \begin{enumerate}
        \item \textbf{Sujeto de estudio:} el propio sistema del PFC en el \textit{compose} de despliegue local.
        \item \textbf{Instrumentos:} Locust para carga, Grafana para observación, scripts de pruebas contractuales y de disponibilidad.
        \item \textbf{Diseño experimental:} bloque completo con $r=10$ repeticiones por escenario; se descartan la primera y la última repetición.
        \item \textbf{Métricas objetivas por característica ISO 25010} según la Tabla~\ref{tab:iso25010}.
        \item \textbf{Análisis:} media, desviación típica e intervalo de confianza al 95\%; se contrasta cada métrica con su umbral objetivo.
        \item \textbf{Amenazas a la validez:} al menos tres internas y dos externas, discutidas en la sección de amenazas del manuscrito.
    \end{enumerate}

    \renewcommand{\arraystretch}{1.2}
    \small
    \begin{xltabular}{\linewidth}{@{}p{3.3cm} X p{2.7cm}@{}}
        \caption{Características ISO/IEC 25010:2023 evaluadas en E4 y sus métricas objetivas.}
        \label{tab:iso25010}\\
        \toprule
        \textbf{Característica} & \textbf{Métrica y protocolo de medición} & \textbf{Umbral objetivo}\\
        \midrule
        \endfirsthead
        \toprule
        \textbf{Característica} & \textbf{Métrica y protocolo de medición} & \textbf{Umbral objetivo}\\
        \midrule
        \endhead
        \bottomrule
        \endfoot
        Fiabilidad & Tasa de errores 5xx durante una hora continua bajo carga nominal. & $<1\%$\\
        Eficiencia de desempeño & Latencia p95 con 50 usuarios concurrentes según escenario Locust. & $<500$ ms\\
        Seguridad & Porcentaje de endpoints con verificación JWT y ausencia de las 10 vulnerabilidades OWASP más comunes~\cite{OWASP2021Top10}. & $100\%$\\
        Mantenibilidad & Cobertura de pruebas unitarias + complejidad ciclomatica media. & $\ge 70\%$; $<10$\\
        Compatibilidad & Aplicación web funciona en Chrome, Firefox y Safari; aplicación móvil funciona en API 26+. & Pasa suite E2E\\
    \end{xltabular}
    \normalsize

\subsection{Módulo H: consolidacion del documento final}

    El manuscrito de la E4 es acumulativo: incluye, actualizadas, todas las secciones de E1, E2 y E3, y añade las secciones nuevas de U5. La cita a las entregas previas es interna: no se referencian como fuentes externas, sino que se integran en un texto continuo con estructura de reporte experimental. El total mínimo es de $20$ páginas de contenido y $\ge 12$ referencias IEEE con DOI o ISBN verificable.

\section{Documento LaTeX exigido}

    La Tabla~\ref{tab:doc-e4} lista la estructura mínima del manuscrito.

    \renewcommand{\arraystretch}{1.15}
    \small
    \begin{xltabular}{\linewidth}{@{}p{3.4cm} X p{1.5cm}@{}}
        \caption{Estructura mínima del manuscrito LaTeX de la Entrega 4.}
        \label{tab:doc-e4}\\
        \toprule
        \textbf{Sección} & \textbf{Contenido} & \textbf{Páginas ref.}\\
        \midrule
        \endfirsthead
        \toprule
        \textbf{Sección} & \textbf{Contenido} & \textbf{Páginas ref.}\\
        \midrule
        \endhead
        \bottomrule
        \endfoot
        Resumen bilingue & Objetivo, método, resultados, conclusiones; $\le 200$ palabras cada uno. & 1\\
        Introducción & Motivación, contribuciones, panorama del semestre. & 1.5\\
        Trabajos relacionados & $\ge 6$ fuentes indexadas sobre BDD, microservicios, móvil y calidad. & 2\\
        Fundamento teórico & Capas + SOLID, GoF, móvil, CI/CD, observabilidad, ISO 25010. & 3\\
        Arquitectura del sistema & Vistas C4 nivel 1 a 3, decisiones (ADR), stack. & 2\\
        Aplicación web & Manual del Módulo B, capturas de las 5 rutas, evidencia i18n. & 2\\
        Aplicación móvil & Manual del Módulo C, capturas, evidencia de capacidades usadas. & 2\\
        Pruebas y CI/CD & Piramide de pruebas, resumen del pipeline, cobertura. & 2\\
        Observabilidad & Módulo F; dashboard como figura; ejemplo de traza distribuida. & 1.5\\
        Evaluación ISO 25010 & Módulo G; tabla de métricas con IC 95\%; discusión. & 2\\
        Discusión y amenazas a la validez & Interpretación, alcance, limitaciones. & 1.5\\
        Conclusiones e trabajo futuro & Cierre del ciclo E1--E4. & 1\\
        Bibliografía & IEEE; $\ge 12$ referencias verificadas. & ---\\
    \end{xltabular}
    \normalsize

\section{Estructura del repositorio}

    El Listado~\ref{lst:repo-e4} muestra la estructura resultante después del \emph{merge} de la E4. Ninguna carpeta de las entregas anteriores desaparece; se añaden las de web, móvil y observabilidad.
    \begin{lstlisting}[language=bash,caption={Estructura del repositorio después de fusionar E4.},label={lst:repo-e4}]
proyecto-pfc/
  README.md                    # actualizado con instrucciones completas
  LICENSE
  .env.example                 # incluye variables de web y móvil
  docker-compose.yml           # backend + CRDB + web + collector
  docs/
    adr/
      ADR-001-arquitectura.md
      ADR-002-microservicios.md
      ADR-003-fragmentación.md
      ADR-004-consenso-raft.md
      ADR-005-patrones-gof.md          # nuevo
      ADR-006-elección-móvil.md        # nuevo
    api/openapi.yaml
    db/schema.sql
    diagrams/
      c4-nivel1.drawio.png
      c4-nivel2.drawio.png
      c4-nivel3.drawio.png             # actualizado con Web + Móvil
    experimentos/
      protocolo-e3.md
      protocolo-e4.md                  # nuevo
      resultados/
        iso25010.csv                   # nuevo
        boxplot_latencia.png           # nuevo
    entrega1/  entrega2/  entrega3/  entrega4/
  services/
    resource-service/                  # refactor en capas
    auth-service/
    notification-service/
  apps/
    web/                               # nuevo, SPA
    mobile/                            # nuevo, Android/multiplataforma
  spark/                               # heredado de E3
  ops/
    grafana/
      pfc-dashboard.json               # nuevo
    prometheus/
      prometheus.yml                   # actualizado
    otel-collector/
      config.yaml                      # nuevo
  release/
    apk/                               # APK generado por CI
    screenshots/                       # capturas de web y móvil
  tests/
    integration/
    load/
    e2e-web/                           # nuevo, Playwright
    contract/                          # nuevo, Pact
  .github/
    workflows/
      ci-cd.yml                        # ampliado con matriz web/móvil
    \end{lstlisting}

\section{Rúbrica de evaluación}

    La rúbrica opera con nueve dimensiones, veintiun criterios y cinco niveles (0 a 4). Se aplica la fórmula
    \begin{equation}
        \text{Nota}_{/10} = \left[ \frac{\sum_i w_i \cdot n_i}{4} \right] \cdot 10 \quad , \quad \sum_i w_i = 1.0
        \label{eq:nota-e4}
    \end{equation}
    Se recuerda que la ausencia total de la aplicación web o de la aplicación móvil aplica un factor multiplicativo de $0{,}5$ sobre la nota final calculada con la Ecuación~\eqref{eq:nota-e4}.

    \renewcommand{\arraystretch}{1.2}
    \small
    \begin{xltabular}{\linewidth}{@{}p{0.9cm} p{4.6cm} X p{1cm}@{}}
        \caption{Rúbrica de evaluación de la Entrega 4.}
        \label{tab:rúbrica-e4}\\
        \toprule
        \textbf{Dim.} & \textbf{Criterio} & \textbf{Indicador clave nivel 4} & \textbf{Peso}\\
        \midrule
        \endfirsthead
        \toprule
        \textbf{Dim.} & \textbf{Criterio} & \textbf{Indicador clave nivel 4} & \textbf{Peso}\\
        \midrule
        \endhead
        \bottomrule
        \endfoot
        D1 & 1.1 Refactor en capas y SOLID & Cuatro capas, dependencias del dominio hacia abstracciones & 6\%\\
           & 1.2 Patrones GoF aplicados & $\ge 5$ patrones justificados en ADR-005; código real, no placebo & 6\%\\
        D2 & 2.1 Aplicación Web (funcional) & Cinco rutas, i18n, tema, JWT correcto, sin errores de consola & 8\%\\
           & 2.2 Aplicación Web (calidad) & Lint verde, cobertura $\ge 70\%$, Dockerfile multi-stage & 4\%\\
        D3 & 3.1 Aplicación Móvil (funcional) & Login, listado y detalle, dos capacidades del dispositivo & 8\%\\
           & 3.2 Aplicación Móvil (calidad) & Tests unit + instrumentados; APK firmado publicado por CI & 4\%\\
        D4 & 4.1 Integración Web + Móvil + Backend & Contratos Pact en verde; misma API sirve a los dos clientes & 5\%\\
           & 4.2 Persistencia distribuida operando & Cluster CRDB de E3 sigue funcionando; healthchecks verdes & 2\%\\
        D5 & 5.1 Piramide de pruebas & Unit + Integración + Contrato + E2E + Carga; artefactos publicados & 6\%\\
           & 5.2 Pipeline CI/CD end-to-end & Los 7 jobs del Listado~\ref{lst:ci} en verde; imagenes en GHCR & 5\%\\
        D6 & 6.1 Observabilidad completa & Métricas + logs con trace\_id + trazas OTel; dashboard versionado & 6\%\\
           & 6.2 Evidencia operativa & Demo en vivo del dashboard bajo carga Locust en la defensa & 2\%\\
        D7 & 7.1 Evaluación ISO 25010 & Cinco características medidas con IC 95\%; discusión honesta & 5\%\\
        D8 & 8.1 Manuscrito final & Estructura de Tabla~\ref{tab:doc-e4}; $\ge 20$ páginas; sin errores LaTeX & 5\%\\
           & 8.2 Bibliografía & IEEE, $\ge 12$ referencias, la mayoría con DOI resoluble & 3\%\\
           & 8.3 Trazabilidad E1--E4 & Se muestra que los artefactos de las entregas previas siguen vivos & 3\%\\
           & 8.4 Repositorio y README & README completo, \path{.env.example} vigente, semillas deterministas & 3\%\\
        D9 & 9.1 Defensa oral (grupal) & Demo en vivo; cada integrante responde preguntas sin leer & 5\%\\
           & 9.2 Reflexión ética & $\ge {}$media página con referencia al ACM Code of Ethics~\cite{ACMCode2018} & 2\%\\
           & 9.3 Discusión de amenazas a la validez & Al menos tres internas y dos externas, cada una con mitigacion & 2\%\\
           & 9.4 Reproducibilidad & \path{docker compose up} + instrucciones despliegan todo desde cero & 2\%\\
    \end{xltabular}
    \normalsize

\section{Anexos}

\subsection*{Anexo A. Lista de verificación final}

    \begin{itemize}
        \item[$\square$] Rama \texttt{feature/entrega-4} fusionada a \texttt{main} por PR con revisión cruzada.
        \item[$\square$] Backend refactorizado en capas; ADR-005 firmado.
        \item[$\square$] Aplicación web (SPA) desplegada con \texttt{docker compose up}; $\ge 5$ rutas; i18n.
        \item[$\square$] Aplicación móvil compilando; APK generado por CI; dos capacidades del dispositivo.
        \item[$\square$] Pipeline CI/CD con los 7 jobs verdes en el último commit de la entrega.
        \item[$\square$] Contratos Pact publicados y verdes; suite E2E Playwright y Espresso pasando.
        \item[$\square$] Prueba de carga Locust ejecutada; capturas incluidas en el manuscrito.
        \item[$\square$] Dashboard Grafana operativo con las seis paneles obligatorios.
        \item[$\square$] Tabla ISO 25010 con las cinco características y sus IC 95\%.
        \item[$\square$] Manuscrito LaTeX compila sin errores; $\ge 20$ páginas de contenido.
        \item[$\square$] Bibliografía con $\ge 12$ referencias IEEE verificadas.
        \item[$\square$] Reflexión ética de al menos media página.
        \item[$\square$] README con instrucciones de despliegue reproducible.
    \end{itemize}

\subsection*{Anexo B. Preguntas obligatorias para la defensa}

    \begin{enumerate}
        \item Justifique la elección de framework para la web y para el móvil, con al menos dos criterios cuantitativos.
        \item Explique como los principios SOLID se materializan en el módulo más complejo de su backend.
        \item Muestre una traza distribuida completa desde la app móvil hasta el cluster CRDB y describa cada \textit{span}.
        \item Explique como el pipeline CI/CD garantiza que el código desplegado cumple los atributos de calidad definidos.
        \item Interprete los resultados de la evaluación ISO 25010: que característica no alcanzo el umbral y por que.
        \item Describa el plan técnico para llevar el sistema a producción real, incluyendo escalado horizontal.
        \item Discuta al menos tres amenazas a la validez de su evaluación experimental y como las mitigo.
        \item Reflexione sobre las responsabilidades eticas del equipo frente a los datos de los usuarios finales.
    \end{enumerate}

\subsection*{Anexo C. Plantilla de ADR (formato Nygard)}

    Cada ADR (\emph{Architectural Decisión Record}) debe seguir el formato:
    \begin{itemize}
        \item \textbf{Título:} número y título corto.
        \item \textbf{Estado:} propuesto, aceptado, sustituido, deprecado.
        \item \textbf{Contexto:} problema, restricciones y alternativas consideradas.
        \item \textbf{Decisión:} elección tomada y razones.
        \item \textbf{Consecuencias:} positivas, negativas y neutrales, con impacto en las entregas anteriores y futuras.
    \end{itemize}

\bibliographystyle{ieeetr}
\bibliography{references_PFC}

\end{document}
